\clearpage\begingroup\SourceLeaf{90b}\HistoricalHead{85}{R}{GAUSS AN ENCKE.}
\begin{center}[9.]\end{center}
\begin{center}\textsc{Gauss} an \textsc{Encke}.\quad Göttingen, 20. August 1833.\end{center}
\noindent . . . . . . . . . . . . . . . . . . . . .\par
\indent Unsere grosse galvanische Kette (6000—7000 Fuss Draht) ist schon lange
ungestört bestehend und schon oft haben wir mit bestem Erfolg ganze kleine
Phrasen einander telegraphisch signalisirt. Ganz besonders merkwürdig und
anfangs für mich überraschend (obwohl es nach richtiger Theorie hätte vorhergesehen
werden können) war der Umstand, dass es dabei gar keiner grossen
Platten oder starken Säuren bedarf; \textls[120]{eine} Kupfer- und Zinkplatte, etwa wie
ein preussischer Thaler gross, und Tuchscheibe mit \textls[120]{reinem Brunnenwasser}
genetzt, ja sogar \textls[120]{destillirtes} Wasser ist vollkommen hinreichend, ja sogar
für die bisherigen Einrichtungen (die ursprünglich nicht \textls[120]{diesen} Zweck hatten)
noch zu stark; die grösste Platte (in sofern man \textls[120]{nur Ein Paar} nimmt) und
starke Säure würde nur etwa eine doppelt so grosse Wirkung geben. . . . . .
\clearpage\endgroup
